
\def \Subject {تعمیم‌پذیری}
\def \Session { یک}
% \setcounter{chapter}{\Session-1}
\renewcommand{\chaptername}{سوال}
\chapter{\Subject}
\section{سوالات توضیحی}
\subsection{مفهوم همبستگی کاذب و شکست \lr{ERM}}
در بسیاری از  الگوریتم‌های مرسوم یادگیری ماشین،‌ تلاش برای کمینه کردن 
ریسک تجربی 
\LTRfootnote{Empirical risk}
است.  از آنجایی که صرف ریسک تجربی مطرح است، ممکن است مدل شروع به یادگیری ویژگی‌های کاذب بکند. این ویژگی‌ها که در داده آموزشی موجود هستند باعث می‌شود تا دقت مدل روی آن‌ها بهبود یابد، ولی از آنجایی که در داده‌های واقعی این ویژگی‌های کاذب وجود ندارند باعث می‌شود تا دقت مدل کاهش یابد.

به عبارت دیگر مدل با یادگیری این ویژگی‌های کاذب دچار بیش‌برازش می‌شود و دیگر به سراغ یادگیری ویژگی‌های اصلی نمی‌رود. 

\subsection{ فرموله سازی خطای بدترین گروه و گروه های داده }

پژوهش‌هایی که برای کاهش سوگیری انجام شده‌اند را می‌توان به دو دسته کلی تقسیم کرد. در دسته اول، ویژگی‌ای که سوگیری را اضافه می‌کند معلوم است. (برای مثال، نمونه عکسی که دسته‌بندی حیوانات را با توجه به پس‌زمینه انجام می‌داد و در مقاله آمده است.) و در دسته دوم این ویژگی‌ها نامعلوم است.

اگر این ويژگی‌ها معلوم نباشند، باید به داده‌ها مراجعه کرد. ویژگی‌های کاذب در بعضی نمونه‌های ورودی بیشتر است. نمونه‌هایی که ویژگی‌های کاذب در آن‌ها بیشتر است، نمونه‌های هم‌راستا با بایاس
\LTRfootnote{Bias-aligned groups}
و نمونه‌هایی که ویژگی‌های کاذب ندارند یا در آن‌ها کمتر است را نمونه‌های در تضاد 
\LTRfootnote{Bias-conflicting groups}
با بایاس
نامیده می‌شوند. 

هدفی که برای مدل تعیین می‌شود این است که مدل روی هر دو این گروه‌ها عملکرد خوبی داشته باشد. مدل‌هایی که دچار سوگیری می‌شوند عملکرد خیلی خوبی روی گروه‌های هم‌راستا با بایاس دارند ولی در گروه‌های در تضاد با بایاس عملکرد ضعیفی خواهند داشت. به همین منظور، هدفی که تعیین می‌شود به این صورت است که مدل باید به شیوه‌ای تنظیم شود که بیشینه خطا در بین گروه‌ها کمینه شود. یعنی اگر مدلی بسیار دچار سوگیری شود، خطای مدل روی گروه‌های در تضاد با بایاس حساب می‌شود.

\begin{equation}
	\min_{\theta} \max_{b \in \mathcal{B}} 
	\left\{
	\hat{\mathcal{L}}_b := \frac{1}{n_b} 
	\sum_{(x,y) \in G_b} \ell(f_\theta(x),y)
	\right\},
	\label{minmax}
\end{equation}

\subsection{نقش تابع زیان \lr{GCE}}
در روش پیشنهادی مقاله، در ابتدا برای اینکه نمونه‌های هم‌راستا با بایاس و در تضاد با بایاس شناسایی شوند، ابتدا یک مدل با سوگیری باید آموزش داده شود تا از عملکرد آن روی داده‌ها بتوان تشخیص داد کدام نمونه‌ها دارای ویژگی‌های کاذب هستند. 

هرچقدر این مدل دارای سوگیری بیشتری باشد، ویژگی‌های کاذب بیشتری شناسایی می‌شوند. پس برای تقویت این موضوع، می‌توان میزان مشتق را برای نمونه‌هایی که مدل درست پیش‌بینی می‌کند بیشتر در نظر گرفت. به این صورت، مدل سعی می‌کند بیشتر بر روی ویژگی‌های کاذب  تکیه کند. زیرا وقتی مقدار مشتق کم باشد، در مرحله بهینه‌سازی، وزن‌ها به مقدار کمی تغییر می‌کنند و با اجرای یک دور آموزش روی همه نمونه‌ها میزان یادگیری مدل برآیندی از ویژگی‌های همه نمونه‌ها خواهد بود ولی در این‌جا مدل بر روی ویژگی‌های خاصی که در بعضی نمونه‌ها وجود دارد تکیه می‌کند. 

برای این‌که مشتق پیش‌بینی‌های درست بیشتر شود از خطای 
\lr{GCE}
استفاده می‌شود که فرمول آن در 
\eqref{gce_formulate}
آمده است. پارامتر 
\lr{q}
مقداری بین صفر تا یک دارد. هرچقدر به یک نزدیک‌تر باشد، تشدید کمتری رخ می‌دهد و هرچقدر مقدار کمتری داشته باشد، باعث می‌شود تا تاثیر نمونه‌هایی که درست دسته‌بندی شده اند بیشتر شود.

\begin{equation}
	\label{gce_formulate}
	\ell = \cfrac{1 - p(y | x)^q}{q}
\end{equation}

\subsection{بررسی فرض اول مقاله}
این فرض بیان می‌کند که مدل دارای سوگیری، دقت بهتری روی نمونه‌های هم‌راستا با بایاس دارد. این فرض از چند جهت مهم است. در مرحله اول با استفاده از این فرض می‌توان فرمول کمینه بیشینه 
\eqref{minmax}
را به این‌صورت ساده کرد که فقط خطا را روی نمونه‌های در تضاد با بایاس حساب نمود. یعنی عملا نیازی به گرفتن بیشینه نیست. (زیرا از قبل می‌دانیم که بیشترین خطا روی نمونه‌های در تضاد با بایاس است.) استفاده دیگری که از این فرض می‌شود، هنگام تشخیص نمونه‌های هم‌راستا با بایاس و در تضاد با بایاس است. از آنجایی که فرض می‌شود خطای مدل روی نمونه‌های در تضاد با بایاس بیشتر است، می‌توان از عکس این فرض استفاده نمود. یعنی نمونه‌هایی که مدل دارای سوگیری روی آن‌ها دقت کمتری دارد، نمونه‌های در تضاد با بایاس هستند.
\subsection{ایده الگوریتم \lr{JTT}}
در این الگوریتم، ابتدا مدل به صورت معمولی آموزش دیده می‌شود. سپس نمونه‌هایی که مدل دقت کمتری روی آن‌ها داشته است‌ (احتمالا نمونه‌های در تضاد با بایاس) وزن بیشتری می‌گیرند و مدل دوباره آموزش می‌بیند. 

این رویکرد به نوعی شبیه به آموزش مدل دارای سوگیری در مقاله قبل است ولی در جهت مخالف. در آن جا با تشدید وزن‌ها سعی می‌شد تا مدل روی نمونه‌های هم‌راستا با سوگیری دقت بهتری داشته باشد و سوگیری بیشتر شود ولی در اینجا با تشدید وزن، سعی می‌شود تا مدل روی نمونه‌های در تضاد با بایاس عملکرد بهتری داشته باشد.
\section{سوالات پیاده‌سازی}
\subsection{آماده‌سازی داده و مصورسازی}
\begin{figure}[H]
	
	\centering
	\includegraphics[width=\textwidth]{images/q1_visualization.jpg}
	\caption{
نمونه داده‌ها پس از تغییرات گفته شده
	}
	\label{q1:visulaization}
\end{figure}
\subsection{آموزش مدل پایه با \lr{ERM}}
\begin{figure}[H]
	
	\centering
	\includegraphics[width=\textwidth]{images/q1_plot_base.jpg}
	\caption{
نمودار خطا و دقت برای مجموعه‌های آموزش و آزمون روی مدل پایه
	}
	\label{q1:plot_base}
\end{figure}
\subsection{استخراج مجموعه خطا (\lr{ERROR SET})}
با مقایسه مقادیر پیش‌بینی شده با مقادیر واقعی،‌ در ۱۰۳۷ مورد، پیش‌بینی با مقدار واقعی یک نیست.   در شکل 
\ref{q1:base_errorset_sample}
نمونه مواردی که غلط دسته‌بندی شده‌اند آمده است.
\begin{figure}[H]
	
	\centering
	\includegraphics[width=\textwidth]{images/q1_base_errorset_sample.jpg}
	\caption{
نمونه موارد خطای مدل پایه روی داده آموزش
	}
	\label{q1:base_errorset_sample}
\end{figure}

\subsection{ساخت دیتاست وزن‌دار \lr{UPSAMPLING}}
بعد از این کار یک مجموعه داده با ۱۱۰۸۱۳ عضو ساخته شد. که این مقدار با محاسبات منطبق است. 
\begin{latin}
	$
	60000 + 1037 * 49 = 110813
	$
\end{latin}
\subsection{آموزش مجدد مدل}
همانطور که در نمودار 
\ref{q1:upscaled_plot}
می‌توان دید، بعد از استفاده از 
\lr{JTT}
مدل توانست به دقت بهتری روی داده‌های آموزش برسد. البته مدل اشباع نشده ولی به دلیل محدودیت تعداد دور آموزش تعیین شده در سوال، از آموزش بیشتر آن صرفنظر شد.
\begin{figure}[H]
	
	\centering
	\includegraphics[width=\textwidth]{images/q1_plot_upscaled.jpg}
	\caption{
نمودارهای خطا و آموزش هنگام آموزش مجدد مدل روی داده‌های 
\lr{upsample}
شده
	}
	\label{q1:upscaled_plot}
\end{figure}
\subsection{پیاده‌سازی تابع زیان \lr{GCE}}
استفاده از تابع زیان 
\lr{GCE}
تاثیر مشهودی دارد
. استفاده از این تابع ضرر باعث می‌شود تا مدل سوگیری بیشتری داشته باشد و حجم 
\lr{errorset}
بیشتر شود. همچنین می‌توان دید که با استفاده از این تابع ضرر، از آنجایی که مدل بیشتر دچار سوگیری می‌شود، دقت روی مجموعه داده آزمون بسیار کاهش می‌یابد. با آموزش مدل با استفاده از 
\lr{GCE}
 حجم 
\lr{errorset}
برابر با ۲۵۴۶ شد که با توجه به سوگیری بیشتر مدل قابل توجیه است. 

با 
\lr{upsample}
کردن این داده‌ها و آموزش مجدد مدل می‌توان به دقتی بهتر از 
اجرای 
\lr{JTT}
با تابع ضرر 
\lr{cross entropy }
رسید که نتایج آن در جدول 
\ref{tab:l1_comp}
آمده است.

\begin{figure}[H]
	
	\centering
	\includegraphics[width=\textwidth]{images/gce_base_plot.jpg}
	\caption{
		نمودارهای خطا و آموزش هنگام آموزش مدل پایه با استفاده از تابع ضرر \lr{GCE} 
	}
	\label{q1:gce_base_plot}
\end{figure}

\begin{figure}[H]
	
	\centering
	\includegraphics[width=\textwidth]{images/q1_loss_acc_gce_upscaled_plots.jpg}
	\caption{
		نمودارهای خطا و آموزش هنگام آموزش مدل 
		\lr{upsample}
		شده
		 با استفاده از تابع ضرر \lr{GCE} 
	}
	\label{q1:gce_upscaled_plot}
\end{figure}
\begin{table}[H]
	\centering
	\caption{مقایسه خطا و دقت نهایی مدل‌ها}
	
	\begin{tabular}{r c c | c c}
		\toprule
		\multirow{2}{*}{مدل} & \multicolumn{2}{c|}{داده‌های آموزش} & \multicolumn{2}{c}{داده‌های آزمون} \\
		\cmidrule(lr){2-3} \cmidrule(lr){4-5}
		& خطا & دقت (\%) & خطا & دقت (\%) \\
		\midrule
		مدل پایه & $0.0825$ & $97.86$ & $0.9778$ & $78.15$ \\
		مدل با \lr{JTT} & $0.0181$ & $99.49$ & $0.7798$ & $84.51$ \\
		مدل با تابع ضرر \lr{GCE} & $0.0614$ & $95.69$ & $1.2067$ & $15.35$ \\
		مدل \lr{jtt} با استفاده از خطاهای \lr{GCE} & $0.0112$ & $99.66$ & $0.7501$ & $88.21$ \\
		\bottomrule
	\end{tabular}
	\label{tab:l1_comp}
\end{table}
